% Abstract (ERL: at most 300 words, no undefined acronyms, no citations; tools/word_count_erl.py
% reports its length). Every number is a \metric{} macro. Revised after the internal review
% (headline on the validated sectors; near-term framing).
\begin{abstract}
\noindent As closed seas and saline lakes shrink, projections of how much seabed will emerge, and of what it will expose, are built on global bathymetry grids. On a flat shelf a small vertical error in such a grid moves the predicted shoreline by kilometres, yet the grids are rarely tested against seabed that has already emerged. The North Caspian Sea, whose running-median level fell by \metric{levels.dahiti.smoothed.range_m}~m between \metric{levels.dahiti.smoothed.date_of_max} and \metric{levels.dahiti.smoothed.date_of_min}, allows that test. We flooded five global grids from the open sea at each level measured by satellite altimetry and compared them with Landsat shorelines of 1992 to 2026. On the bare western and interfluve coasts, where the observed waterlines are independently validated, \metric{rev.s.item1.shape.validated4.n_products_all} of the grids misplace them by \metric{rev.w.h1.coasts3.d_km_min}--\metric{rev.w.h1.coasts3.d_km_max}~km even after an area-fitted vertical shift, under every check against whole-metre depths, coarse cells and level error; on the northeastern coast for most grids and checks. On the open shelf the error cannot be told apart from those effects. The shift matches their median flooded areas to the observed, but not across levels or out of sample. The fifth grid, ETOPO 2022, fits above \metric{h1.disp_band.edge_2_m}~m because its shoreline there comes from land heights acquired after emergence and from a lake outline of 2000, never from lake depths. At Great Salt Lake, Utah, where none of the grids holds lake-bed soundings, all five fail the same test (one passes one arm under a pooled floor), while a lidar-based elevation model passes it. Measured from the observed 2026 shoreline, the grids place \metric{rev.n.composite.summary.common_location.1.min_km2}--\metric{rev.n.composite.summary.common_location.1.max_km2}~km\textsuperscript{2} of seabed in the first metre of further fall, or \metric{rev.n.summary.shifted.1.min_km2}--\metric{rev.n.summary.shifted.1.max_km2}~km\textsuperscript{2} once their offsets are removed. The water-level pathway decides how much seabed emerges by mid-century; the grids decide where.
\end{abstract}
